\chapterNote{Based on a talk given on December 22, 2022}

\chapter{A Sublime Abiding}
\tocChapterNote{Based on a talk given on December 22, 2022}

Those of us who have taken monastic vows and those who are committed lay practitioners receive much encouragement to keep in mind the long-term potential of practice: the ultimate goal of realizing \mbox{liberation.} The Buddha talks about it in many ways in the teachings. He instructs us not to become complacent in \mbox{the practice.}

The Dhamma teacher Ruth Denison, who passed away a number of years ago, emphasized the same idea towards the end of her life. A group of people went down from Abhayagiri not long before she passed away to visit her and to wish her well. One of her notable messages was: ``Don't settle for mere happiness. You have the possibility for full liberation.''

The encouragement for realizing the complete goal is a focus of our teachings. We gear our attention towards those teachings that help lead us to the experience of the Unborn, the Uncreated, the Unoriginated, and the Unformed. We keep returning to that goal.

At the same time, we live in the world and have to negotiate our ways around the world, whether it's in the monastery or outside of the monastery. We try to make it a useful place to practice and create skillful, supportive conditions. We have to learn how to negotiate the conditioned realm as part of our \linebreak path to the Unconditioned.

\looseness=-1
We are nearing the end of December, and the days are getting darker and shorter. The nights are long. There is a coolness in the air. The signals of nature are moving us towards quiet, turning inwards, contemplation, \mbox{and stillness.}

It's delightful to experience that at this time of year. The holiday messages at their best, such as those around Christmas, honor the birth and life of a very compassionate figure. The uplifting side of Christmas is remembering qualities of nurturing, care, concern, compassion, and skillfulness in our relationships. There is more attention on family, friends, and the people with whom we live, and on the goodness that can be cultivated from these relationships.

It's also easy to get lost in some of the less desirable aspects of the season. There is the side of consumer mentality and acquisitiveness, the material world. Businesses make a lot of money at this time of year and can get caught up in seeking gain by selling their products; others can get very absorbed in the consumer ethic.

Going into a shopping mall at this time of year, with gaudy decorations and blaring music, can be jolting. For those of us who are living in or around the monastery and lead a quiet life, it's an affront to the senses. It's a mixed time of year, depending on what we focus on.

Learning how to live skillfully in the conditioned realm is part of the path to the Unconditioned. This means learning how to live in the world, but not be of the world. We learn how not to get lost in the worldliness of the world and how to navigate in a way that is still conducive to the path.

Even in the monastery, we don't isolate ourselves and shun the world. We do try to step away from the worldliness of unskillful actions and nurture skillful ways of being in the world, but we're not apart from it. One of the most skillful ways of developing ourselves in the realm of the conditioned is the development and abiding in the four \emph{brahmavihāras}, the divine abidings.

\emph{Vihāra} means dwelling place, so brahmavihāra means dwelling or abiding in the mind state of a brahma god. These are beautiful states of mind that beings in high spiritual realms dwell in for long periods of time. We can also cultivate them and make them an abiding in our lives in this world. Even though these states are of the conditioned realm, they are a very important part of the path to total freedom.

The four brahmavihāras are \emph{mettā} (loving-kindness), \emph{karuṇā} (compassion), \emph{muditā} (sympathetic joy), and \emph{upekkhā} (equanimity). Many of the current teachings regarding the brahmavihāras focus on specific meditation techniques to develop them. The emphasis is on an active practice of \emph{mettā-bhāvanā} using various mental recitations, verbal phrases, such as, ``May I be well and happy.'' There certainly is value in this approach for many people. For others, including myself, an alternative approach to developing the divine abidings is first investigating what mind state might be obstructing the natural arising of the brahmavihāras. We look closely at the \emph{kilesas}, the obstructive mind states that stand in the way of opening to sublime mind states.

If we don't clearly see the mental obstructions that block the brahmavihāras, then we're possibly only doing it halfway. If we don't notice the underlying tendencies that lurk as potentials, and if we don't see them clearly, then it's so easy for them to catch us off guard, even if we don't think we have such states of mind very often. If we don't attend to them with clarity and honesty, then they may rise up unexpectedly, and take us by surprise. As with the First Noble Truth, we have to fully understand the experience of \emph{dukkha} before we can start to let it go.

Obstructive mind states are stimulated by outside circumstances and we have to meet them with honesty and clear awareness, without getting lost in them and without moving into and indulging in them.

\looseness=-1
For the first brahmavihāra of mettā, the obstructive states that can rear up are ill will, anger, aversion, and all of their unfortnate manifestations. Somebody who is unskillful or unpleasant to be around may trigger something in us that rises up as ill will, aversion, or anger. We shouldn't develop aversion to the aversion and try to turn away from it, pretending it's not there. We shouldn't bury it or feel like we shouldn't feel that way. We should turn gently towards the aversion to fully understand it, to fully face it and see that state of mind as it is experienced right in the present moment. This is done most skillfully by attending to the bodily energies.

If we get involved in the storyline too much, ``He did this. She said that,'' then we're just increasing the proliferative tendency. We need to drop the story and feel the feeling. We need to turn to the affective tone of the experience to really understand it without getting lost in proliferation. We should develop the bravery to face the unpleasant feeling. There may be clenching in the stomach, tightness in the cheeks, contraction of the facial muscles, or the rising of heat. Sometimes, \linebreak if the feeling is quite intense, there's a flush that's really unpleasant, and our normal reaction is to try and get away from it. Nobody wants to have unpleasant feelings. In order not to deal with the pain, we can get lost in the storyline, trying to justify, explain, or present our counterarguments, with all kinds of conversations \mbox{in our head.}

With patience, though, we can get past that discomfort. We can bear up and experience the pain by facing it and allowing it to manifest as an energetic quality through the body. We can gently bring our breathing, kind attention, and spaciousness to the experience and not get wrapped up in the storyline. We can attend to the contracted state on the bodily level and then more deeply in the heart. We allow ourselves the space to be fully with the experience, without absorbing into and encouraging it.

Fully attending to intense bodily energies takes a lot of patience and spacious kindness and further encourages the qualities of kindness, spaciousness, gentleness, and warmth to arise. It doesn't mean that we're endorsing unskillful behavior, but we're learning to experience a broader sense of kindness, including well-wishing to ourselves, by owning up to the truth that people make mistakes. We make mistakes and other people make mistakes. We can disentangle ourselves from the enmeshed obstructive mind state and wish ourselves and others freedom from those reactions.

\enlargethispage{1\baselineskip}
We can realize that people who are triggering our aversion, anger, and ill will, are probably feeling the same. We can wish them well. If they can experience relaxation from tense mind states, then they will act in ways that are much more pleasant for us, too. It's a win-win situation.

The afflictive state of wishing harm is the opposite of karuṇā or compassion, the second brahmavihāra. We might see this arise in subtle ways when somebody who has been unskillful receives the karmic consequence of their action and experiences some misfortune as a result. When the tables turn and they get their comeuppance, we might feel a momentary, brief delight that they are ``getting theirs'' now. This happens a lot in the world of politics, where there is so much anger, tension, and contention. A wish for revenge emerges when the tides turn, which they always do.

\looseness=-1
There may be some insensitivity or cruelty when somebody is experiencing suffering and difficulty. There may be an inclination, for whatever reason, to turn a cold shoulder to that person. That's a really unpleasant state of mind. If that happens, even on a small scale, what does that feel like? We need to turn our attention towards the unpleasant feeling of constriction and tightness when we experience moments of cool indifference or insensitivity towards somebody else's suffering.

By experiencing the results of our own insensitivity, we realize that other people are no different than we are. Compassion means learning to be with the passions and difficulties as they manifest and recognizing that somebody else's suffering is the same as our suffering.

\enlargethispage{1\baselineskip}
Compassion is the quality of quivering or trembling with, to be in resonance with someone else's suffering. When we know our own suffering, then the natural response to another's suffering will be to feel it at some level. As much as we want to be free from suffering, we want others to be free from suffering. It's all the same.

The third brahmavihāra is sympathetic joy. An afflictive quality that we can experience related to the third brahmavihāra is envy when someone else is experiencing success. Envy is an unpleasant state of mind that arises when we want to get someone's attention, care, or concern, but they are giving it to somebody else. As they are giving someone support or offering others praise and recognition, we're thinking, ``Well, what about me? What about \emph{me}?'' We ignore the fact that somebody else is receiving justifiable recognition or support and that they're feeling good. We get lost in our ``what about me?''

\looseness=-1
When that envy arises, we can respond with the same approach as with aversion. We fully experience the bodily energy and say to ourselves, ``Okay, well, this is some sort of conditioned wish for recognition.'' As long as we are subject to praise and blame and success and failure, when somebody else gets something we would like, envy can arise. We need to know this honestly and not get lost in self-flagellation for experiencing afflictive states, because they are part of our conditioning, part of the patterns that we have developed for whatever reason. We need to recognize them as energetic states and then allow them to pass. If we attend to the energetic state with kind attention, it will dissipate.

Once the negative state passes, we can turn towards that other being who is receiving something pleasant and recognize that when we receive that kind of attention, we also feel pleasure. To resonate with another person's success is a rare commodity in the world. To really open up and feel genuine delight when someone else has received something is quite wonderful. That is muditā, the brahmavihāra of sympathetic joy.

The last brahmavihāra is equanimity. When we're not equanimous we are pulled around, as with the eight worldly winds. We get pulled towards pleasure and away from pain. We are buffeted around by seeking gain and avoiding loss. We are affected by praise and blame and by seeking good reputation and avoiding disrepute. We are consumed with bouncing back and forth in reaction to the ways of the world, the world of opposites. Realizing the way the world is, we can choose either to get lost in it, repeatedly being overwhelmed by opposing states of mind, or we can choose to see them with clarity. We can develop the ability to be with these opposites with a clear, knowing state of mind. Equanimity has a very strong quality of discernment, which we can employ in our encounters with the world. We can learn to step back from overreacting to conditions because they are always going to fluctuate. They \mbox{will always change.}

Our basic practice mode is to own up to the underlying states that obstruct the brahmavihāras. We then encourage the states of kindness, sensitivity, \mbox{compassion}, warmth, well-wishing, and quiet joy. We experience our own gains and successes and those of others, without attachment. We can experience evenness and an understanding of the ways of the world. These qualities are dwelling places in our hearts where we can settle, not just as a practice, but as an actual abiding. We steadily make that our frame of reference. By seeing the disadvantage of how we normally operate, we step away from the normal ways of the world. We then can live in the heavenly realms in our minds, allowing that space to open up and develop.

\looseness=-1
Equanimity is so important in and of itself, and it's an important part of the other three brahmavihāras, too. When we have equanimity, we are less likely to get lost in what the commentaries call the near enemies of the brahmavihāras. Equanimity provides an alert clarity and awareness of our states of mind. In regards to the development of mettā, for example, equanimity prevents us from sliding into attachment, which is the near enemy of mettā. With clarity, we move beyond \linebreak the personal fondness for one specific person and cultivate a global quality of kindness and well-wishing, independent of individuality. Mettā is a feeling for all beings and a very sublime state to abide in. Attachment narrows our attention down to one individual or a limited number of people that we happen to like or \mbox{are attracted to.}

When we are developing loving-kindness reflections, we shouldn't let our minds slide into affection for an individual or individuals. Even if we are not moving into sexual desire, Māra is oftentimes waiting on the edges for an opportunity to cloud the mind. The wonderful quality of loving-kindness for all beings can turn into something more involved with attachment, clinging, and confusion. Equanimity and wisdom help to prevent us from sliding into the near enemy of mettā.

The same is true with karuṇā, compassion. When we are feeling compassion for someone and their suffering, if we are not careful, not mindful, not equanimous, we can slide into a feeling of overwhelm. We become identified with the suffering of someone else, and we don't really know how to handle it. We absorb into it and don't have the proper boundaries that come from clarity. Whether it's our own or someone else's, we get overwhelmed with the suffering of the world. Equanimity keeps compassion stable, clear, and balanced, so we know when and how we can help without getting lost in emotion or reactivity.

The brahmavihāra of muditā can tend towards over- \linebreak excitement, even to the point of agitation about \mbox{something} good. When a family member, partner, or monastic friend is experiencing something wonderful and we get too excited, the sense of balance disappears. We've lost the perspective of ``well, yeah, this is good and to be delighted in, but it's also part of the \mbox{conditioned realm.''}

\looseness=-1
To keep even a good thing going takes constant management, and things will probably shift. One may not always experience this kind of success, and the tides might turn. The blame might happen again. The \linebreak loss might happen again. This isn't meant to negate the good quality of the experience or to look at it as sour grapes: ``Oh, it's going to go away anyway.'' We can delight in goodness \emph{and} also reflect that it's a conditioned experience that will change and fade over time.

Equanimity enhances and stabilizes the other three brahmavihāras. However, equanimity itself can sink into its near enemy of indifference. Sometimes we can slide into thinking, ``Oh, well, everything's impermanent, so what does it matter?'' If everything is impermanent and it's going to change, does that lead us to indifference or even nihilism? If we let equanimity slide in that direction, we lose the benefits of all of the other brahmavihāras. We won't be able to sustain a sense of well-wishing (the wish for the end of suffering) or delight in joyful experiences.

Indifference is not at all what equanimity is. Equanimity involves clear seeing. The states of mind that we are trying to develop are the results of intentional skillful reflection and good actions. They are the results of skillful kamma. Equanimity includes the wisdom that knows how these states come into being and how we can maintain them as skillful qualities. This is Right Effort: raising up, establishing, and maintaining the wholesome qualities that serve as parts of the path.

As we develop the four brahmavihāras, we also realize that we can develop them to the point where they are ``abundant, exalted, and immeasurable.'' These are very expansive, spacious states of mind. By letting go of constricted states of mind, the mind opens, relaxes, expands, and broadens. The experience of each of the individual four brahmavihāras grows into a very wide, open state of awareness.

But it all comes back to recognizing and owning up to the \emph{kilesas}, the constricted states of mind, and letting them go. This is done through clear understanding, not by avoiding them, and by allowing ourselves to experience them. As we let go of these underlying states, we can open up to the comfort and \mbox{happiness that arises.}

The brahmavihāras are the best we can do in the con-ditioned realm. Even on our quest towards realizing the Unconditioned, we have to negotiate our way through the world and react in the most skillful way to the conditions around us. The more we do that, the more this light ease of being can manifest. The old habits slowly wear away. While we are still part of the world, we develop a more beautiful way to be in the world. The world becomes a place where we can comfortably abide as we negotiate closer to realizing the Unconditioned.\hfill\linebreak
